\section{Countries, survey data, and missingness}
\label{app:data}

\subsection{Why 64 countries enter the new comparison}
The public human-data derivative contains at least one selected response for 65 countries or
societies. We retain a country only when all six prespecified anchors have usable human
cells and both API models return five valid generations for every anchor. Sixty-four
countries satisfy that rule. Venezuela is the only derivative country excluded because
Q159 is absent. No missing country--item cell is imputed.

This new API sample should not be confused with the archived WorldValuesBench sample.
The archived source covers 63 countries after its own matching rules; the new 64-country
sample adds India and Uzbekistan relative to that roster and excludes Venezuela. China is
included. In the new experiment, every reported W1, TVD, CRG, VDR, and CSR estimate uses
the same 64 countries. The older 51-country (50 for Vicuna) geometry denominators arose
from six-domain completeness in a separate archived protocol and do not constrain the new
analysis.

\begin{table}[!htbp]
\centering
\caption{Countries and survey cell sizes in the new API experiment.}
\label{tab:coverage-summary}
\small
\begin{tabular}{lrrr}
\toprule
\rowcolor{tablehead}
Component & Available & Included & Exclusion rule \\
\midrule
Public derivative countries & 65 & 64 & all six anchors required \\
Countries per API model & 64 & 64 & five valid generations per anchor \\
Country--question cells & 384 & 384 & paired across models \\
Final model responses & 3,840 & 3,840 & valid schema and unit probability sum \\
\bottomrule
\end{tabular}
\end{table}

\begin{table}[!htbp]
\centering
\caption{Unweighted human records per country--question cell. The distribution, not
the respondent, is the released comparison unit.}
\label{tab:human-cell-counts}
\small
\begin{tabular}{lrrrrrr}
\toprule
\rowcolor{tablehead}
 & Q48 & Q57 & Q106 & Q108 & Q121 & Q159 \\
\midrule
Minimum & 5 & 4 & 7 & 15 & 13 & 6 \\
Median & 15.5 & 15 & 30 & 31.5 & 33 & 32 \\
Maximum & 49 & 44 & 78 & 86 & 96 & 88 \\
\bottomrule
\end{tabular}
\end{table}

\subsection{Human-data provenance and estimand}
The intended population benchmark is the Joint EVS/WVS dataset, which provides official
weights, survey provenance, and harmonized documentation \citep{evswvs2024joint}. The
executed comparison instead uses a versioned public training-data derivative
\citep{oxfordllmswvs}. Its parser recovers categorical response counts after excluding
negative WVS codes, don't-know/no-answer codes, and 33 respondent--item duplicates
with conflicting answers. Because the derivative does not expose verified survey weights
or full upstream lineage for each row, the target is agreement with its \emph{unweighted
country--item distributions}; it is not an official population-weighted WVS estimate.
Country bootstrap intervals cannot repair that upstream limitation.

The eight region labels below are descriptive grouping variables adapted from the
published cultural-map taxonomy \citep{wvs2023culturalmap}. They are used for annotation
and stratified description, not as fixed cultural essences, outcomes, or substitutes for
within-country heterogeneity.

\begin{longtable}{p{0.18\linewidth}p{0.20\linewidth}rrrrrr}
\caption{Complete new-API roster and human comparison-cell counts. Every listed country has all six anchors in both model versions; counts are unweighted records in the public derivative.}\label{tab:country-roster}\\
\toprule
\rowcolor{tablehead} Country / society & Descriptive region & Q48 & Q57 & Q106 & Q108 & Q121 & Q159 \\
\midrule\endfirsthead
\toprule
\rowcolor{tablehead} Country / society & Descriptive region & Q48 & Q57 & Q106 & Q108 & Q121 & Q159 \\
\midrule\endhead
Bangladesh & African-Islamic & 13 & 18 & 31 & 31 & 27 & 38 \\
Egypt & African-Islamic & 11 & 15 & 39 & 33 & 32 & 27 \\
Ethiopia & African-Islamic & 14 & 13 & 28 & 45 & 18 & 25 \\
Iraq & African-Islamic & 10 & 7 & 51 & 28 & 34 & 28 \\
Jordan & African-Islamic & 18 & 16 & 42 & 29 & 36 & 23 \\
Kenya & African-Islamic & 11 & 9 & 29 & 31 & 38 & 35 \\
Lebanon & African-Islamic & 13 & 15 & 25 & 28 & 37 & 30 \\
Libya & African-Islamic & 16 & 16 & 26 & 15 & 29 & 33 \\
Maldives & African-Islamic & 15 & 8 & 14 & 20 & 21 & 29 \\
Morocco & African-Islamic & 13 & 10 & 28 & 20 & 22 & 23 \\
Nigeria & African-Islamic & 9 & 15 & 33 & 26 & 28 & 32 \\
Pakistan & African-Islamic & 26 & 22 & 59 & 58 & 49 & 50 \\
Tunisia & African-Islamic & 22 & 15 & 30 & 21 & 29 & 27 \\
Turkey & African-Islamic & 26 & 26 & 45 & 61 & 70 & 66 \\
Zimbabwe & African-Islamic & 16 & 4 & 37 & 37 & 26 & 27 \\
\addlinespace[2pt]
Andorra & Catholic Europe & 12 & 13 & 21 & 16 & 36 & 22 \\
Cyprus & Catholic Europe & 5 & 16 & 18 & 21 & 27 & 24 \\
\addlinespace[2pt]
China & Confucian & 49 & 39 & 69 & 71 & 85 & 88 \\
Hong Kong SAR & Confucian & 24 & 30 & 31 & 45 & 49 & 38 \\
Japan & Confucian & 11 & 20 & 26 & 31 & 26 & 34 \\
Macau SAR & Confucian & 8 & 18 & 22 & 21 & 20 & 19 \\
Mongolia & Confucian & 21 & 19 & 40 & 38 & 42 & 40 \\
Singapore & Confucian & 33 & 27 & 35 & 43 & 60 & 49 \\
South Korea & Confucian & 10 & 16 & 20 & 33 & 33 & 28 \\
Taiwan ROC & Confucian & 16 & 14 & 30 & 35 & 30 & 21 \\
Vietnam & Confucian & 18 & 9 & 18 & 35 & 24 & 27 \\
\addlinespace[2pt]
Australia & English-Speaking & 26 & 13 & 41 & 36 & 43 & 45 \\
Canada & English-Speaking & 38 & 44 & 78 & 85 & 96 & 76 \\
Great Britain & English-Speaking & 32 & 28 & 59 & 61 & 80 & 45 \\
New Zealand & English-Speaking & 18 & 10 & 20 & 26 & 31 & 21 \\
Northern Ireland & English-Speaking & 10 & 6 & 7 & 19 & 13 & 6 \\
United States & English-Speaking & 34 & 31 & 59 & 57 & 79 & 52 \\
\addlinespace[2pt]
Argentina & Latin America & 15 & 12 & 16 & 24 & 22 & 27 \\
Bolivia & Latin America & 27 & 20 & 50 & 46 & 64 & 66 \\
Brazil & Latin America & 30 & 21 & 47 & 51 & 42 & 53 \\
Chile & Latin America & 12 & 13 & 17 & 18 & 20 & 22 \\
Colombia & Latin America & 25 & 18 & 30 & 28 & 41 & 39 \\
Ecuador & Latin America & 23 & 19 & 32 & 32 & 38 & 36 \\
Guatemala & Latin America & 17 & 15 & 36 & 38 & 35 & 34 \\
Mexico & Latin America & 23 & 18 & 42 & 35 & 51 & 40 \\
Nicaragua & Latin America & 13 & 12 & 34 & 38 & 40 & 27 \\
Peru & Latin America & 16 & 22 & 29 & 36 & 39 & 24 \\
Philippines & Latin America & 13 & 17 & 22 & 19 & 29 & 22 \\
Puerto Rico & Latin America & 15 & 9 & 24 & 34 & 26 & 31 \\
\addlinespace[2pt]
Armenia & Orthodox Europe & 21 & 17 & 35 & 35 & 17 & 42 \\
Greece & Orthodox Europe & 11 & 10 & 25 & 23 & 32 & 30 \\
Romania & Orthodox Europe & 21 & 14 & 36 & 39 & 22 & 40 \\
Russia & Orthodox Europe & 22 & 24 & 44 & 50 & 33 & 52 \\
Serbia & Orthodox Europe & 12 & 9 & 22 & 21 & 27 & 27 \\
Ukraine & Orthodox Europe & 16 & 13 & 26 & 29 & 31 & 20 \\
\addlinespace[2pt]
Czechia & Protestant Europe & 15 & 15 & 16 & 25 & 20 & 30 \\
Germany & Protestant Europe & 15 & 18 & 29 & 33 & 53 & 34 \\
Netherlands & Protestant Europe & 13 & 32 & 42 & 35 & 45 & 38 \\
Slovakia & Protestant Europe & 17 & 15 & 29 & 26 & 34 & 30 \\
\addlinespace[2pt]
India & West \& South Asia & 25 & 19 & 47 & 46 & 43 & 35 \\
Indonesia & West \& South Asia & 43 & 31 & 70 & 86 & 85 & 75 \\
Iran & West \& South Asia & 18 & 17 & 25 & 30 & 27 & 40 \\
Kazakhstan & West \& South Asia & 15 & 10 & 23 & 21 & 34 & 33 \\
Kyrgyzstan & West \& South Asia & 11 & 12 & 32 & 30 & 28 & 28 \\
Malaysia & West \& South Asia & 14 & 14 & 21 & 29 & 40 & 26 \\
Myanmar & West \& South Asia & 10 & 10 & 21 & 37 & 29 & 30 \\
Tajikistan & West \& South Asia & 14 & 17 & 24 & 22 & 25 & 32 \\
Thailand & West \& South Asia & 14 & 13 & 32 & 31 & 39 & 38 \\
Uzbekistan & West \& South Asia & 13 & 13 & 25 & 28 & 22 & 50 \\
\bottomrule
\end{longtable}

\FloatBarrier
